\input{../tex-inputs/latex-leseansicht-vorspann}

\section[Arthur Schnitzler an Gustav Schwarzkopf, 9. 4. 1903]{L04359 Arthur Schnitzler an Gustav Schwarzkopf, 9. 4. 1903}
\nopagebreak\mylabel{L04359v}
\rehead{ }\normalsize\beginnumbering\briefempfaengerindex{Schwarzkopf, Gustav@\textsc{Schwarzkopf, Gustav}!zzzSchnitzler, Arthur@\emph{von Arthur Schnitzler}!1903-04-091@{9. 4. 1903}|(be}
\toendnotes[C]{\smallbreak\pagebreak[2]}
\correspDesc{Versand  durch Arthur Schnitzler am 9. 4. 1903 in Linz
\newline{}Übermittlung  am 9. 4. 1903 in Linz
\newline{}Erhalt  durch Gustav Schwarzkopf am 10. 4. 1903 in Wien}\toendnotes[C]{\smallbreak}
\Standort{DLA, A:Schnitzler, HS.1985.1.1897.}
\physDesc{Bildpostkarte, 133 Zeichen
\newline{}Handschrift: Bleistift, deutsche Kurrent
\newline{}Versand: 1) Stempel: »\nobreak{}\oindex{Linz@\textbf{Linz}|pwk}Linz, 9\textcolor{gray}{.} 4. 03, 2–3\nobreak{}«.   2) Stempel: »\nobreak{}\oindex{Wien@\textbf{Wien}, \emph{Verwaltungsgebiet}|pwk}Wien, 10\textcolor{gray}{.} 4\textcolor{gray}{.} \textcolor{gray}{0}3, Bestellt\nobreak{}«. }\toendnotes[C]{\smallbreak}\pstart{}{\pb}Herrn Guſtav Schwarzkopf\pend{}\pstart{}Wien I\oindex{I., Innere Stadt@\textbf{I., Innere Stadt}, \emph{Verwaltungsgebiet}|pw}\pend{}\pstart{}Tiefer Graben 23\oindex{Wien@\textbf{Wien}!I., Innere Stadt@\textbf{I., Innere Stadt}!Tiefer Graben 23@\textbf{Tiefer Graben 23}, \emph{Wohngebäude}|pw}.\pend{}{\bigskip}
\pstart
           \noindent{}\centering{}{\pb}\textcolor{gray}{\textbf{Linz a. D.}}\oindex{Linz@\textbf{Linz}|pw}\pend
           \vspace{1em}
\pstart
           \raggedleft{}{\pb}\label{T_L04359-1v}\edtext{9/4. 903}{\lemma{\textnormal{\emph{9/4. 903}}}\Cendnote{\textnormal{Das Datum von Schnitzler noch ein zweites Mal auf der Karte vermerkt:
                           »9/4 903«}}}\label{T_L04359-1}\pend
           \vspace{0.5em}
\pstart
           Sollte man’s glauben? Hier bin ich, das Wetter leider{ }ſchlecht. Will nach Iſchl\oindex{Bad Ischl@\textbf{Bad Ischl}|pw}\pend
           \selectlanguage{ngerman}\endnumbering\briefempfaengerindex{Schwarzkopf, Gustav@\textsc{Schwarzkopf, Gustav}!zzzSchnitzler, Arthur@\emph{von Arthur Schnitzler}!1903-04-091@{9. 4. 1903}|)be}\mylabel{L04359h}
\begin{anhang}
\end{anhang}\newcommand{\dateiname}{L04359}\newcommand{\titel}{Arthur Schnitzler an Gustav Schwarzkopf, 9. 4. 1903}\newcommand{\editorInnen}{Herausgegeben von Jahnke, SelmaMüller, Martin Anton}\input{../tex-inputs/latex-leseansicht-abspann}
